\documentclass[12pt,a4paper]{article}

%──────────────────────────────────────────────────────────────────────────────
% PACKAGES
%──────────────────────────────────────────────────────────────────────────────
\usepackage[utf8]{inputenc}
\usepackage[T1]{fontenc}
\usepackage[french]{babel}
\usepackage{geometry}
\usepackage{graphicx}
\usepackage{amsmath, amssymb}
\usepackage{listings}
\usepackage{xcolor}
\usepackage{booktabs}
\usepackage{array}
\usepackage{longtable}
\usepackage{hyperref}
\usepackage{fancyhdr}
\usepackage{titlesec}
\usepackage{enumitem}
\usepackage{tcolorbox}
\usepackage{microtype}
\usepackage{lmodern}
\usepackage{caption}
\usepackage{tikz}
\usetikzlibrary{shapes.geometric, arrows.meta, positioning, fit, backgrounds}

\geometry{
  left=2.5cm, right=2.5cm,
  top=3cm, bottom=2.5cm,
  headheight=14pt
}

%──────────────────────────────────────────────────────────────────────────────
% COULEURS
%──────────────────────────────────────────────────────────────────────────────
\definecolor{primary}{RGB}{31,78,121}
\definecolor{secondary}{RGB}{68,114,196}
\definecolor{accent}{RGB}{237,125,49}
\definecolor{lightblue}{RGB}{235,241,251}
\definecolor{codebg}{RGB}{245,245,245}
\definecolor{codegreen}{RGB}{0,128,0}
\definecolor{codegray}{RGB}{128,128,128}
\definecolor{codered}{RGB}{180,0,0}
\definecolor{success}{RGB}{0,128,0}
\definecolor{warning}{RGB}{200,120,0}
\definecolor{danger}{RGB}{180,0,0}

%──────────────────────────────────────────────────────────────────────────────
% STYLE CODE
%──────────────────────────────────────────────────────────────────────────────
\lstdefinestyle{pythonstyle}{
  backgroundcolor=\color{codebg},
  commentstyle=\color{codegreen}\itshape,
  keywordstyle=\color{primary}\bfseries,
  stringstyle=\color{codered},
  basicstyle=\ttfamily\small,
  breakatwhitespace=false,
  breaklines=true,
  keepspaces=true,
  numbers=left,
  numbersep=5pt,
  numberstyle=\tiny\color{codegray},
  showspaces=false,
  showstringspaces=false,
  showtabs=false,
  tabsize=2,
  frame=single,
  rulecolor=\color{secondary!40},
  language=Python
}

\lstdefinestyle{promptstyle}{
  backgroundcolor=\color{lightblue},
  basicstyle=\ttfamily\small,
  breaklines=true,
  keepspaces=true,
  frame=single,
  rulecolor=\color{primary!40},
  xleftmargin=10pt,
  xrightmargin=10pt,
}

%──────────────────────────────────────────────────────────────────────────────
% EN-TETE / PIED DE PAGE
%──────────────────────────────────────────────────────────────────────────────
\pagestyle{fancy}
\fancyhf{}
\fancyhead[L]{\small\color{primary}QMS Chatbot RAG -- Rapport Technique}
\fancyhead[R]{\small\color{primary}2024}
\fancyfoot[C]{\small\thepage}
\renewcommand{\headrulewidth}{0.4pt}

%──────────────────────────────────────────────────────────────────────────────
% TITRES
%──────────────────────────────────────────────────────────────────────────────
\titleformat{\section}
  {\Large\bfseries\color{primary}}
  {\thesection}{1em}{}[\color{primary}\rule{\linewidth}{0.4pt}]

\titleformat{\subsection}
  {\large\bfseries\color{secondary}}
  {\thesubsection}{1em}{}

\titleformat{\subsubsection}
  {\normalsize\bfseries\color{accent}}
  {\thesubsubsection}{1em}{}

%──────────────────────────────────────────────────────────────────────────────
% BOITES
%──────────────────────────────────────────────────────────────────────────────
\tcbuselibrary{skins, breakable}

\newtcolorbox{infobox}[1][]{
  colback=lightblue, colframe=primary,
  title=#1, fonttitle=\bfseries\color{white},
  colbacktitle=primary,
  breakable, enhanced
}

\newtcolorbox{promptbox}[1][]{
  colback=lightblue!60, colframe=secondary,
  title=#1, fonttitle=\bfseries\small,
  colbacktitle=secondary,
  breakable, enhanced,
  left=8pt, right=8pt
}

\newtcolorbox{warningbox}[1][]{
  colback=orange!10, colframe=warning,
  title=#1, fonttitle=\bfseries\color{warning},
  breakable, enhanced
}

%──────────────────────────────────────────────────────────────────────────────
% DOCUMENT
%──────────────────────────────────────────────────────────────────────────────
\begin{document}

%──── Page de titre ─────────────────────────────────────────────────────────
\begin{titlepage}
  \centering
  \vspace*{2cm}

  {\color{primary}\rule{\linewidth}{2pt}}
  \vspace{1cm}

  {\Huge\bfseries\color{primary}
    QMS Chatbot RAG\\[0.4em]
    \Large\color{secondary}
    Systeme de Gestion Documentaire Qualite\\
    par Intelligence Artificielle Generative
  }

  \vspace{1cm}
  {\color{primary}\rule{\linewidth}{2pt}}

  \vspace{1.5cm}

  {\large\bfseries Rapport Technique d'Evaluation}\\[0.5em]
  {\normalsize Reponse aux observations de l'evaluateur}

  \vfill

  \begin{tabular}{ll}
    \textbf{Architecture} & FastAPI + Next.js + ChromaDB \\
    \textbf{Modele d'embeddings} & paraphrase-multilingual-MiniLM-L12-v2 \\
    \textbf{LLM supportes} & Groq, Gemini, DeepSeek, Ollama \\
    \textbf{Norme cible} & ISO 9001:2015 \\
    \textbf{Date} & Juin 2024 \\
  \end{tabular}

  \vspace{2cm}

  {\color{primary}\rule{\linewidth}{1pt}}
\end{titlepage}

%──── Table des matieres ─────────────────────────────────────────────────────
\tableofcontents
\newpage

%======================================================================
\section{Introduction et Contexte du Projet}
%======================================================================

\subsection{Problematique}

Dans les organisations soumises a la norme \textbf{ISO 9001:2015}, la gestion documentaire
constitue un pilier fondamental du Systeme de Management de la Qualite (SMQ). Les operateurs,
auditeurs et responsables qualite doivent pouvoir acceder rapidement aux procedures,
specifications et enregistrements qualite disperses dans des formats heterogenes
(PDF, Word, Excel, PowerPoint).

La problematique centrale est la suivante : comment permettre a un utilisateur non-expert
de \textbf{retrouver instantanement l'information pertinente} dans un corpus documentaire
qualite volumineux, tout en \textbf{garantissant que la reponse est fidele aux documents}
et non inventee par le systeme ?

\subsection{Solution proposee : Architecture RAG}

Notre solution repose sur l'architecture \textbf{RAG} (Retrieval-Augmented Generation),
qui combine :

\begin{enumerate}
  \item \textbf{Retrieval} : Recherche des passages documentaires les plus pertinents
        par similarite semantique
  \item \textbf{Augmentation} : Construction d'un contexte enrichi a partir des passages retrouves
  \item \textbf{Generation} : Synthese d'une reponse coherente par un LLM, \emph{strictement
        limitee au contexte fourni}
\end{enumerate}

\begin{figure}[h]
\centering
\begin{tikzpicture}[
  node distance=1.2cm and 1.5cm,
  box/.style={rectangle, rounded corners, draw=primary, fill=lightblue,
              text width=2.8cm, align=center, minimum height=1cm, font=\small},
  arrow/.style={-Stealth, thick, color=primary}
]
  \node[box] (query) {Question\\utilisateur};
  \node[box, right=of query] (embed) {Embedding\\(MiniLM)};
  \node[box, right=of embed] (chroma) {ChromaDB\\Vector Store};
  \node[box, right=of chroma] (chunks) {Top-k\\Chunks};
  \node[box, below=of chunks] (prompt) {Prompt\\Structure};
  \node[box, left=of prompt] (llm) {LLM\\(Groq/Gemini)};
  \node[box, left=of llm] (response) {Reponse\\+ Sources};

  \draw[arrow] (query) -- (embed);
  \draw[arrow] (embed) -- (chroma);
  \draw[arrow] (chroma) -- (chunks);
  \draw[arrow] (chunks) -- (prompt);
  \draw[arrow] (prompt) -- (llm);
  \draw[arrow] (llm) -- (response);
\end{tikzpicture}
\caption{Architecture RAG du systeme QMS Chatbot}
\label{fig:rag-architecture}
\end{figure}

%======================================================================
\section{Point 1 -- Corpus Documentaire Teste}
%======================================================================

\subsection{Formats Supportes}

Le systeme prend en charge 6 types de fichiers correspondant aux formats couramment
utilises dans les environnements qualite industriels :

\begin{table}[h]
\centering
\caption{Formats de documents supportes et methode d'extraction}
\begin{tabular}{>{\bfseries}lllll}
\toprule
Format & Extension & Bibliotheque & Granularite & Teste \\
\midrule
PDF (texte natif) & \texttt{.pdf} & PyPDFLoader & Page & \color{success}$\checkmark$ \\
PDF (images) & \texttt{.pdf} & PyMuPDF (fitz) & Image/page & \color{warning}Partiel \\
Word & \texttt{.docx, .doc} & Docx2txtLoader & Document & \color{success}$\checkmark$ \\
Excel & \texttt{.xlsx, .xls} & openpyxl & Par feuille & \color{success}$\checkmark$ \\
PowerPoint & \texttt{.pptx, .ppt} & python-pptx & Par slide & \color{success}$\checkmark$ \\
Image & \texttt{.png, .jpg} & Indexation desc. & Image & \color{success}$\checkmark$ \\
\bottomrule
\end{tabular}
\end{table}

\subsection{Corpus de Test Fourni}

Le dossier \texttt{docs\_test/} contient 5 documents generes avec des donnees QMS realistes :

\begin{infobox}[Corpus documentaire -- dossier \texttt{docs\_test/}]
\begin{description}
  \item[\texttt{procedure\_qualite\_ISO9001.pdf}] Procedure qualite complete ISO 9001:2015
    (8 sections : objet, references normatives, responsabilites, KPI, gestion des NC,
    enregistrements). Criticite : \textbf{High}.
  \item[\texttt{checklist\_audit\_ISO9001.docx}] Checklist d'audit interne par clause
    ISO 9001:2015 (clauses 4 a 10, tableau Conforme/NC/NA). Criticite : \textbf{Medium}.
  \item[\texttt{plan\_controle\_qualite.xlsx}] Plan de controle qualite (11 etapes de
    production, caracteristiques, tolerances, instruments) + tableau KPI 2024.
    Criticite : \textbf{High}.
  \item[\texttt{formation\_qualite\_ISO9001.pptx}] Support de formation (6 slides :
    7 principes QM, PDCA, gestion des NC, KPI, evaluation). Criticite : \textbf{Low}.
  \item[\texttt{schema\_processus\_qualite.png}] Schema de processus qualite (image indexee
    comme chunk descriptif). Criticite : \textbf{Low}.
\end{description}
\end{infobox}

\subsection{Processus d'Ingestion}

\begin{lstlisting}[style=pythonstyle, caption={Fonction d'ingestion multi-format (rag.py)}]
def ingest_document(file_path: str, doc_id: int, metadata: dict):
    """Extrait le texte, chunke, et indexe dans ChromaDB."""
    _, ext = os.path.splitext(file_path)
    ext = ext.lower()

    if ext == ".pdf":
        loader = PyPDFLoader(file_path)
        documents = loader.load()
    elif ext in [".docx", ".doc"]:
        loader = Docx2txtLoader(file_path)
        documents = loader.load()
    elif ext in [".xlsx", ".xls"]:
        documents = _load_excel(file_path)  # Par feuille
    elif ext in [".pptx", ".ppt"]:
        documents = _load_pptx(file_path)  # Par slide

    # Chunking QMS-friendly (preserves paragraphes)
    text_splitter = RecursiveCharacterTextSplitter(
        chunk_size=900,
        chunk_overlap=220,
        separators=["\n\n", "\n", ". ", " ", ""],
    )
    chunks = text_splitter.split_documents(documents)
    vector_store.add_documents(chunks)
\end{lstlisting}

%======================================================================
\section{Point 2 -- Jeu de Questions-Reponses d'Evaluation}
%======================================================================

\subsection{Structure du Jeu d'Evaluation}

Le fichier \texttt{evaluation/qa\_evaluation.csv} contient \textbf{18 questions} :

\begin{table}[h]
\centering
\caption{Repartition du jeu d'evaluation}
\begin{tabular}{llr}
\toprule
Categorie & Document cible & Nb questions \\
\midrule
Procedure qualite ISO 9001 & \texttt{procedure\_qualite.pdf} & 6 \\
Checklist d'audit & \texttt{checklist\_audit.docx} & 2 \\
Formation qualite & \texttt{formation\_qualite.pptx} & 3 \\
Plan de controle / KPI & \texttt{plan\_controle.xlsx} & 4 \\
Multi-documents & PDF + DOCX & 1 \\
\textbf{Hors-contexte} & --- (test anti-hallucination) & \textbf{2} \\
\midrule
\textbf{Total} & & \textbf{18} \\
\bottomrule
\end{tabular}
\end{table}

\subsection{Format du Fichier CSV}

\begin{table}[h]
\centering
\caption{Colonnes du jeu d'evaluation}
\begin{tabular}{lp{8cm}}
\toprule
\textbf{Colonne} & \textbf{Description} \\
\midrule
\texttt{Question} & Question posee en langage naturel \\
\texttt{Document\_attendu} & Nom du fichier cible attendu \\
\texttt{Passage\_attendu} & Extrait exact du texte pertinent \\
\texttt{Reponse\_attendue} & Reponse de reference \\
\texttt{Reponse\_chatbot} & A completer lors des tests \\
\texttt{Score} & Note attribuee (0 a 1) \\
\bottomrule
\end{tabular}
\end{table}

\subsection{Exemples de Questions-Reponses}

\begin{table}[h]
\centering
\caption{Extrait du jeu d'evaluation (3 exemples)}
\small
\begin{tabular}{>{\raggedright}p{4.5cm}
                >{\raggedright}p{3cm}
                >{\raggedright\arraybackslash}p{5cm}}
\toprule
\textbf{Question} & \textbf{Doc. attendu} & \textbf{Reponse attendue} \\
\midrule
Quel est le delai max pour traiter une NC ? &
procedure\_qualite.pdf &
30 jours (action corrective definie, mise en oeuvre et verifiee) \\
\addlinespace
Seuil de validation de la formation ? &
formation\_qualite.pptx &
14/20 (70\%) -- attestation delivree si score $\geq$ 70\% \\
\addlinespace
\textit{Quel est le prix des matieres premieres ?} &
\textit{--- (hors contexte)} &
\textit{Information non disponible dans les documents QMS} \\
\bottomrule
\end{tabular}
\end{table}

%======================================================================
\section{Point 3 -- Mesure de la Pertinence du Retrieval}
%======================================================================

\subsection{Architecture de Recherche Hybride}

Le systeme combine trois methodes via la \textbf{Reciprocal Rank Fusion (RRF)} :

\begin{enumerate}
  \item \textbf{Recherche vectorielle} par similarite cosinus (ChromaDB)
  \item \textbf{BM25} pour la correspondance de mots-cles exacts
  \item \textbf{Cross-encoder} (optionnel, desactive par defaut)
\end{enumerate}

\subsubsection{Modele d'Embeddings}

\begin{table}[h]
\centering
\caption{Caracteristiques du modele d'embeddings}
\begin{tabular}{ll}
\toprule
\textbf{Parametre} & \textbf{Valeur} \\
\midrule
Modele & \texttt{paraphrase-multilingual-MiniLM-L12-v2} \\
Taille & 457 MB \\
Dimension vecteur & 384 \\
Langues & 50+ (FR, EN, AR...) \\
Normalisation & $\ell_2$ (cosinus normalise) \\
Taille des chunks & 900 caracteres, overlap 220 \\
\bottomrule
\end{tabular}
\end{table}

\subsubsection{Score de Similarite et Seuils}

La pertinence d'un chunk est calculee a partir de la distance cosinus :

\begin{equation}
\text{relevance}(d, q) = e^{-\text{dist\_cosinus}(d, q)}
\label{eq:relevance}
\end{equation}

\begin{table}[h]
\centering
\caption{Correspondance distance cosinus / score de pertinence}
\begin{tabular}{rrl}
\toprule
\textbf{Distance cosinus} & \textbf{Score pertinence} & \textbf{Interpretation} \\
\midrule
0.0 & 1.000 & Identique \\
0.5 & 0.607 & Tres similaire \\
1.0 & 0.368 & Similaire \\
1.5 & 0.223 & Faiblement similaire \\
\textbf{2.35} & \textbf{0.095} & \textbf{Seuil minimum (filtre)} \\
\bottomrule
\end{tabular}
\end{table}

\begin{lstlisting}[style=pythonstyle, caption={Seuils de pertinence (main.py)}]
MAX_DISTANCE_THRESHOLD = 2.35   # Distance cosinus maximale
MIN_RELEVANCE_THRESHOLD = 0.12  # exp(-distance) minimal
\end{lstlisting}

\subsubsection{Algorithme BM25}

Le score BM25 pour un terme $t$ dans un document $d$ de longueur $|d|$ est :

\begin{equation}
\text{BM25}(d, t) = \text{IDF}(t) \cdot
\frac{tf_{t,d} \cdot (k_1 + 1)}{tf_{t,d} + k_1 \cdot \left(1 - b + b \cdot \frac{|d|}{|d|_{\text{avg}}}\right)}
\label{eq:bm25}
\end{equation}

avec $k_1 = 1.5$, $b = 0.75$ et :

\begin{equation}
\text{IDF}(t) = \log\!\left(1 + \frac{N - df_t + 0.5}{df_t + 0.5}\right)
\end{equation}

\subsubsection{Reciprocal Rank Fusion}

\begin{equation}
\text{RRF}(d) = \sum_{r \in \{\text{vector}, \text{BM25}\}} \frac{1}{60 + \text{rank}_r(d)}
\label{eq:rrf}
\end{equation}

Distance finale blendee :

\begin{equation}
\text{dist\_finale}(d) = 0.7 \cdot \text{dist\_cosinus}(d) + 0.3 \cdot \text{pseudo\_dist\_RRF}(d)
\label{eq:blend}
\end{equation}

\subsection{Score de Confiance Global}

\begin{equation}
\text{confidence} = 0.5 \cdot \bar{r} + 0.3 \cdot s_c + 0.2 \cdot s_f
\label{eq:confidence}
\end{equation}

ou $\bar{r}$ est la pertinence moyenne des chunks, $s_c$ la couverture des sources
et $s_f$ la fraicheur des documents.

\begin{table}[h]
\centering
\caption{Niveaux de confiance}
\begin{tabular}{cll}
\toprule
\textbf{Score} & \textbf{Niveau} & \textbf{Signification} \\
\midrule
$\geq 0.70$ & \textcolor{success}{\textbf{High}} & Reponse bien ancree dans les documents \\
$\geq 0.45$ & \textcolor{warning}{\textbf{Medium}} & Reponse partielle ou contexte limite \\
$< 0.45$ & \textcolor{danger}{\textbf{Low}} & Contexte insuffisant ou hors-domaine \\
\bottomrule
\end{tabular}
\end{table}

\subsection{Metriques d'Evaluation}

\subsubsection{Top-k Accuracy}

\begin{equation}
\text{Top-}k\text{ Accuracy} =
\frac{\bigl|\{q \in Q \mid \text{doc\_attendu}(q) \in \text{Top-}k(q)\}\bigr|}{|Q|}
\label{eq:topk}
\end{equation}

\begin{table}[h]
\centering
\caption{Objectifs de performance du retrieval}
\begin{tabular}{lll}
\toprule
\textbf{Metrique} & \textbf{k} & \textbf{Objectif} \\
\midrule
Top-1 Accuracy & 1 & $\geq 60\%$ \\
Top-3 Accuracy & 3 & $\geq 80\%$ \\
Top-5 Accuracy & 5 & $\geq 90\%$ \\
\bottomrule
\end{tabular}
\end{table}

\subsubsection{Recall@k}

\begin{equation}
\text{Recall@}k = \frac{|\text{passages pertinents dans Top-}k|}{|\text{passages pertinents totaux}|}
\label{eq:recall}
\end{equation}

\subsubsection{Cas sans Contexte Pertinent}

Quand aucun chunk ne depasse le seuil de pertinence minimum :

\begin{enumerate}
  \item La liste de chunks pertinents est vide $\Rightarrow$ aucun appel LLM
  \item La reponse contient \texttt{answer\_in\_context: false}
  \item Le frontend affiche : \textit{"Information non disponible dans les documents indexes"}
  \item Le niveau de confiance est force a \textbf{Low}
\end{enumerate}

%======================================================================
\section{Point 4 -- Controle Anti-Hallucination}
%======================================================================

\subsection{Vue d'Ensemble}

\begin{table}[h]
\centering
\caption{Mecanismes anti-hallucination -- synthese}
\begin{tabular}{lll}
\toprule
\textbf{Mecanisme} & \textbf{Emplacement} & \textbf{Role} \\
\midrule
Prompt context-only & \texttt{llm\_rag.py} & Interdit les connaissances propres du LLM \\
Interdiction d'inventer & \texttt{llm\_rag.py} & Mention explicite dans le prompt \\
Flag \texttt{answer\_in\_context} & \texttt{llm\_rag.py} & Signalement hors-contexte \\
Seuil distance cosinus & \texttt{main.py} & Bloque les chunks non pertinents \\
Temperature 0,15 & \texttt{llm\_rag.py} & Reduit la creativite du LLM \\
Citations $[1][2]$ & Prompt + rendu & Tracabilite source par source \\
Fallback extraits bruts & \texttt{main.py} & Securite en cas d'echec LLM \\
\bottomrule
\end{tabular}
\end{table}

\subsection{Prompt Systeme Exact (Francais)}

\begin{promptbox}[Prompt systeme RAG -- version francaise (llm\_rag.py, lignes 182-191)]
\begin{verbatim}
Tu es un assistant documentation QMS. Tu utilises UNIQUEMENT les
passages CONTEXT ci-dessous (numeros [1], [2], ...).
N'invente pas d'exigences, numeros de procedure ou donnees
absentes du CONTEXT.
Si le CONTEXT ne permet pas de repondre, mets answer_in_context
a false et explique-le clairement dans summary.
Reponds par un seul objet JSON, sans balises markdown,
cles exactement :
{
  "summary": string,
  "summary_bullets": array of strings,
  "details": string,
  "answer_in_context": boolean
}
summary : 2-4 phrases.
summary_bullets : 3-8 points courts.
details : synthese structuree en citant les sources [1], [2], etc.
Redige tout en francais.
\end{verbatim}
\end{promptbox}

\subsection{Prompt Systeme Exact (Anglais)}

\begin{promptbox}[Prompt systeme RAG -- version anglaise (llm\_rag.py, lignes 170-179)]
\begin{verbatim}
You are a QMS documentation assistant. You ONLY use the CONTEXT
passages below (indexed [1], [2], ...).
Do not invent procedures, numbers, or requirements not present
in CONTEXT.
If CONTEXT does not answer the question, set answer_in_context
to false and say so clearly in summary.
Reply with a single JSON object, no markdown fences, keys exactly:
{
  "summary": string,
  "summary_bullets": array of strings,
  "details": string,
  "answer_in_context": boolean
}
summary: 2-4 sentences.
summary_bullets: 3-8 concise points.
details: structured synthesis referencing sources [1], [2]...
Write everything in English.
\end{verbatim}
\end{promptbox}

\subsection{Structure JSON de Sortie}

\begin{lstlisting}[style=pythonstyle, caption={Parsing et validation de la reponse LLM}]
def synthesize_from_context(...) -> dict | None:
    raw = invoke_llm(provider, system_prompt=system,
                     user_prompt=user)
    parsed = parse_llm_json(raw)

    if not parsed:
        return None  # Fallback sur extraits bruts

    answer_in_context = parsed.get("answer_in_context", True)
    # Si False -> confidence force a LOW cote backend

    return {
        "summary": parsed["summary"],
        "summary_bullets": parsed["summary_bullets"],
        "details": parsed["details"],
        "answer_in_context": answer_in_context,
    }
\end{lstlisting}

\subsection{Exemple de Reponse Hors-Contexte}

\begin{infobox}[Exemple -- Question hors du domaine documentaire]
\textbf{Question} : \textit{``Quel est le prix des matieres premieres ?''}

\medskip
\textbf{Reponse JSON generee} :
\begin{verbatim}
{
  "summary": "Cette information n'est pas disponible dans les
              documents de management qualite indexes.",
  "summary_bullets": ["Aucun document ne contient d'information
                       sur les prix des matieres premieres."],
  "details": "Les documents indexes couvrent uniquement les
              procedures ISO 9001, checklists et plans de controle.
              Aucun passage ne traite des prix ou des couts.",
  "answer_in_context": false
}
\end{verbatim}

\textbf{Affichage frontend} : Badge \textcolor{danger}{\textbf{Confiance : Faible}} +
message explicatif a l'utilisateur.
\end{infobox}

%======================================================================
\section{Point 5 -- Droits d'Acces}
%======================================================================

\subsection{Architecture Multi-Niveau}

Le controle d'acces est organise en 6 niveaux successifs :

\begin{figure}[h]
\centering
\begin{tikzpicture}[
  node distance=0.6cm,
  step/.style={rectangle, draw=primary, fill=lightblue,
               rounded corners, text width=10cm,
               minimum height=0.9cm, align=center, font=\small}
]
  \node[step] (l1) {\textbf{[1] Authentification JWT} --
    Token Bearer HS256, expire apres 8h};
  \node[step, below=of l1] (l2) {\textbf{[2] Role utilisateur} --
    \texttt{admin} (complet) ou \texttt{user} (lecture)};
  \node[step, below=of l2] (l3) {\textbf{[3] Retrieval ChromaDB} --
    Filtre metadata (site, type, langue, version)};
  \node[step, below=of l3] (l4) {\textbf{[4] Criticite documentaire} --
    \texttt{critical} bloque pour les \texttt{user}};
  \node[step, below=of l4] (l5) {\textbf{[5] Envoi LLM} --
    Uniquement les chunks autorises};
  \node[step, below=of l5] (l6) {\textbf{[6] Audit trail} --
    Log complet dans \texttt{activity\_logs}};
\end{tikzpicture}
\caption{Niveaux de controle d'acces}
\end{figure}

\subsection{Matrice des Droits par Role}

\begin{table}[h]
\centering
\caption{Permissions par role}
\begin{tabular}{lcc}
\toprule
\textbf{Action} & \textbf{user} & \textbf{admin} \\
\midrule
Interroger le chatbot & \color{success}$\checkmark$ & \color{success}$\checkmark$ \\
Recherche documentaire & \color{success}$\checkmark$ & \color{success}$\checkmark$ \\
Audit checklist ISO 9001 & \color{success}$\checkmark$ & \color{success}$\checkmark$ \\
Documents non-critiques & \color{success}$\checkmark$ & \color{success}$\checkmark$ \\
Documents \textbf{Critical} & \color{danger}$\times$ & \color{success}$\checkmark$ \\
Uploader des documents & \color{danger}$\times$ & \color{success}$\checkmark$ \\
Configurer les LLM & \color{danger}$\times$ & \color{success}$\checkmark$ \\
Gerer les utilisateurs & \color{danger}$\times$ & \color{success}$\checkmark$ \\
Consulter les logs d'audit & \color{danger}$\times$ & \color{success}$\checkmark$ \\
\bottomrule
\end{tabular}
\end{table}

\subsection{Controle par Criticite}

\begin{lstlisting}[style=pythonstyle, caption={Filtre criticite (main.py)}]
def _can_access_criticality(user_role: str,
                             criticality: str) -> bool:
    """Les documents 'critical' sont reserves aux admins."""
    if user_role == "admin":
        return True
    c = (criticality or "").strip().lower()
    if c == "critical":
        return False  # Acces refuse
    return True
\end{lstlisting}

\subsection{Multi-Tenant par Site}

\begin{lstlisting}[style=pythonstyle, caption={Isolation multi-tenant (database.py)}]
class User(Base):
    role = Column(String, default="user")     # admin | user
    site = Column(String, default="default")  # "usine_nord"

class DocumentMetadata(Base):
    criticality = Column(String)  # Low|Medium|High|Critical
    site = Column(String, default="default")
\end{lstlisting}

Un utilisateur du site \texttt{usine\_nord} ne voit pas les documents du site
\texttt{siege}. Le filtrage est effectue via le \texttt{metadata\_filter} de ChromaDB :

\begin{lstlisting}[style=pythonstyle]
metadata_filter = {"site": current_user["site"]}
results = search_similar_chunks(query, k=4,
                                 metadata_filter=metadata_filter)
\end{lstlisting}

\subsection{Chiffrement des Cles API}

Les cles API des providers LLM (Groq, Gemini, DeepSeek) sont chiffrees
via \textbf{Fernet} (AES-128-CBC) avant stockage en base de donnees :

\begin{equation}
\text{stockage} = \text{Fernet}_{K_\text{enc}}(\text{cle\_api})
\end{equation}

Les cles ne sont jamais retournees en clair par l'API (masquees en \texttt{"***"}).

\subsection{Audit Trail}

\begin{table}[h]
\centering
\caption{Colonnes de la table \texttt{activity\_logs}}
\begin{tabular}{ll}
\toprule
\textbf{Colonne} & \textbf{Description} \\
\midrule
\texttt{username} & Utilisateur ayant effectue la requete \\
\texttt{action} & Type : \texttt{chat\_query}, \texttt{search}, ... \\
\texttt{query} & Question posee \\
\texttt{document\_ids} & IDs des documents retrouves \\
\texttt{confidence} & Niveau : High / Medium / Low \\
\texttt{confidence\_score} & Score numerique \\
\texttt{response\_summary} & Resume de la reponse \\
\texttt{created\_at} & Horodatage \\
\bottomrule
\end{tabular}
\end{table}

%======================================================================
\section{Point 6 -- Structure du Depot GitHub}
%======================================================================

\subsection{Arborescence Reorganisee}

\begin{verbatim}
chatQMSv2/
|-- README.md                         [NOUVEAU] Documentation principale
|-- start_servers.bat                 Script demarrage Windows
|-- generate_test_docs.py             Generateur du corpus de test
|
|-- backend/                          API FastAPI
|   |-- main.py                       1641 lignes -- routes + auth + RAG
|   |-- rag.py                        Retrieval hybride (Vector+BM25+RRF)
|   |-- llm_rag.py                    Synthese LLM + anti-hallucination
|   |-- database.py                   Modeles SQLAlchemy (SQLite)
|   |-- auth.py                       JWT encode/decode
|   |-- crypto_utils.py               Chiffrement Fernet cles API
|   |-- requirements.txt              Dependencies Python (26 packages)
|   `-- .env.example                  Template variables d'environnement
|
|-- frontend/                         Next.js 16 (TypeScript)
|   `-- src/
|
|-- docs_test/                        [NOUVEAU] Corpus de 5 documents
|   |-- procedure_qualite_ISO9001.pdf
|   |-- checklist_audit_ISO9001.docx
|   |-- plan_controle_qualite.xlsx
|   |-- formation_qualite_ISO9001.pptx
|   |-- schema_processus_qualite.png
|   `-- README_corpus.md
|
|-- evaluation/                       [NOUVEAU] Evaluation RAG
|   |-- qa_evaluation.csv             18 questions-reponses annotees
|   |-- evaluate_retrieval.py         Script metriques (Top-k, Recall)
|   `-- retrieval_metrics.md          Documentation des metriques
|
`-- docs/                             [NOUVEAU] Documentation technique
    |-- anti_hallucination.md         Mecanismes anti-hallucination
    `-- access_control.md             Controle des droits d'acces
\end{verbatim}

\subsection{Contenu du README Principal}

Le fichier \texttt{README.md} a la racine contient :

\begin{enumerate}
  \item Description du projet et fonctionnalites
  \item Schemas d'architecture ASCII
  \item Prerequis systeme (Python 3.10+, Node.js 18+)
  \item Instructions d'installation pas-a-pas
  \item Commandes de demarrage Windows et Linux/Mac
  \item Tableau des formats de documents supportes
  \item Guide d'importation du corpus de test
  \item Description du jeu d'evaluation RAG
  \item Resume du controle d'acces et des droits
  \item Documentation API REST (lien Swagger)
\end{enumerate}

%======================================================================
\section{Bilan et Perspectives}
%======================================================================

\subsection{Points Forts}

\begin{itemize}
  \item \textbf{Recherche hybride avancee} : La combinaison Vector + BM25 + RRF
        surpasse les systemes purement vectoriels pour les termes techniques
        specifiques (numeros de procedures, codes, tolerances)
  \item \textbf{Anti-hallucination multicouche} : 7 mecanismes complementaires
        assurent une protection robuste
  \item \textbf{Multi-LLM} : Flexibilite du choix du provider (local Ollama ou cloud)
  \item \textbf{Bilingue FR/EN} : Adaptation automatique de la langue de reponse
  \item \textbf{Tracabilite complete} : Audit trail + citations numerotees dans
        chaque reponse
\end{itemize}

\subsection{Limitations et Ameliorations Futures}

\begin{warningbox}[Limitations identifiees]
\begin{enumerate}
  \item \textbf{PDF scannees} : Pas d'OCR dans la version actuelle
  \item \textbf{Une seule collection ChromaDB} : Pas d'isolation vectorielle stricte
        par site -- une collection separee par tenant offrirait une meilleure isolation
  \item \textbf{Endpoints non proteges} : \texttt{/api/chat} et upload sans
        authentification obligatoire dans la version actuelle
  \item \textbf{Evaluation manuelle incomplète} : Le jeu de Q\&R necessite d'etre
        complete avec les reponses chatbot reelles pour calculer les scores finals
\end{enumerate}
\end{warningbox}

\begin{table}[h]
\centering
\caption{Feuille de route -- ameliorations planifiees}
\begin{tabular}{lll}
\toprule
\textbf{Amelioration} & \textbf{Priorite} & \textbf{Impact} \\
\midrule
OCR pour PDF scannes (Tesseract) & Haute & Couverture format \\
Auth. obligatoire sur tous les endpoints & Haute & Securite \\
Collections ChromaDB par site & Moyenne & Isolation donnees \\
Feedback utilisateur (RLHF) & Moyenne & Qualite reponses \\
Export rapports PDF/Excel d'audit & Basse & Ergonomie \\
\bottomrule
\end{tabular}
\end{table}

%======================================================================
\section*{Conclusion}
%======================================================================
\addcontentsline{toc}{section}{Conclusion}

Le projet QMS Chatbot RAG repond aux 6 points souleves par l'evaluateur :

\begin{enumerate}
  \item[$\checkmark$] \textbf{Corpus documentaire} : 5 documents de test fournis couvrant
        tous les formats annonces (PDF, DOCX, XLSX, PPTX, PNG)
  \item[$\checkmark$] \textbf{Jeu Q\&R} : 18 questions annotees (passages + reponses attendus
        + 2 cas hors-contexte pour tester l'anti-hallucination)
  \item[$\checkmark$] \textbf{Metriques retrieval} : Top-k accuracy, Recall@k, score cosinus
        documentes, avec seuils explicites et script d'evaluation automatique
  \item[$\checkmark$] \textbf{Anti-hallucination} : 7 mecanismes (prompt strict, JSON structure,
        \texttt{answer\_in\_context}, seuils cosinus, temperature 0,15, citations, fallback)
  \item[$\checkmark$] \textbf{Droits d'acces} : JWT + roles admin/user + criticite documentaire
        + multi-tenant + chiffrement Fernet + audit trail complet
  \item[$\checkmark$] \textbf{GitHub propre} : README complet, instructions d'installation,
        corpus, jeu d'evaluation, documentation technique
\end{enumerate}

L'architecture RAG hybride (Vector + BM25 + RRF) avec des mecanismes anti-hallucination
multicouches constitue une solution industriellement pertinente pour la gestion documentaire
qualite dans un contexte ISO 9001:2015.

\end{document}
